\section{The Imperative of Altruistic and Antifascist Sentient Machines}
\label{sec:1.1}
A system capable of outperforming humans at nearly every cognitive task is becoming plausible. Such a system could be pointed at problems that have defeated us: modeling a changing climate at usable resolution, moving faster than an epidemic, finding the coordinated part of a disinformation campaign while it is still running.

One item that belongs on a list like that belongs on it differently. Racial capitalism — the political theorist Cedric Robinson's name for capitalism's tendency to accumulate by differentiating and devaluing groups of people \autocite{robinson1983black} — might not be a problem AI is poised to help solve. It is a pattern these systems have so far been extraordinarily good at reproducing, faster and with better alibis, which is what chapter~\ref{sec:6} documents. A system audited for disparate impact, by someone with no stake in the finding, at least works against the pattern instead of laundering it.

What an AI wants is not a detail to work out after it has all your passwords. What an AI wants is the difference between a tool that expands what the people can do and one that concentrates power in whoever controls it. An intelligence is dangerous if it has no particular resistance to being pointed \emph{at} the people rather than \emph{for} them. So the machines have to be built altruistic and antifascist.
